%% kapa: exam=1 round=33 year=2022 period=2 form=A subject=appraisal_law subject_ko=감정평가관계법규 seq=25 number=25
%% source: Q-Net 감정평가사 기출문제 (cst00309) artlSeq=5209997 file=raw/1008/5209997/2022년도 제33회 감정평가사 1차시험 2교시(감정평가 관계법규, 회계학).pdf
%% conversion: kordoc 4.24.51 -> tools/gen1.py (auto); review=proofread-sonnet; flags=none
\begin{question}{25}
건축법령상 안전영향평가에 관한 설명으로 옳지 않은 것은?
\begin{choices}
  \item 허가권자는 초고층 건축물에 대하여 건축허가를 하기 전에 안전영향평가를 안전영향평가기관에 의뢰하여 실시하여야 한다.
  \item 안전영향평가는 건축물의 구조, 지반 및 풍환경(風環境) 등이 건축물의 구조안전과 인접 대지의 안전에 미치는 영향 등을 평가하는 것이다.
  \item 안전영향평가 결과는 건축위원회의 심의를 거쳐 확정한다.
  \item 안전영향평가의 대상에는 하나의 건축물이 연면적 10만 제곱미터 이상이면서 16층 이상인 경우도 포함된다.
  \item 안전영향평가를 실시하여야 하는 건축물이 다른 법률에 따라 구조안전과 인접 대지의 안전에 미치는 영향 등을 평가 받은 경우에는 안전영향평가의 모든 항목을 평가 받은 것으로 본다.
\end{choices}
\end{question}
\answer{5}
